\chapter{رفلاکس وزیکویورترال}

\VURpage{۱}{20261005\_024348639.jpg}
\section*{\lr{VUR}}
\begin{itemize}
\item بازگشت ادرار از مثانه به حالب و کلیه.
\item زمانی رخ می‌دهد که تونل ساب‌موکوزال بین مخاط و عضلهٔ دترسور کوتاه باشد یا کلاً وجود نداشته باشد.
\item ۱–۲\% کودکان.
\item معمولاً مادرزادی و در اغلب موارد فامیلی است.
\item در ۳۰\% کودکان دختر با \lr{UTI} وجود دارد.
\item در ۵–۱۵\% شیرخواران با هیدرونفروز آنتی‌ناتال.
\end{itemize}
\lr{VUR} یک ریسک فاکتور مهم برای پیلونفریت است.

\lr{VUR} عامل \lr{reflux nephropathy} یا همان \lr{reflux-related renal injury}.

در کودک با \lr{UTI} تب‌دار اگر \lr{VUR} وجود داشته باشد، احتمال آسیب کلیه ۳ برابر است.

اگر آسیب کلیه در زمینهٔ \lr{VUR} شدید باشد: \lr{renin-mediated HTN}.

جنین پیش از تولد ممکن است آسیب کلیه پیدا کند یا حتی بمیرد.

\section*{دسته‌بندی \lr{VUR}}
\begin{itemize}
\item بر اساس \lr{VCUG} (سیستویورتروگرام حین ادرار): \lr{I--V}.
\item \lr{VUR}: پرایمری؛ \lr{Secondary}.
\item اختلال عملکرد مثانه–روده (\lr{BBD}) می‌تواند \lr{VUR} را تشدید کند.
\end{itemize}
\subsection*{سندرم مگاسیستیس مگایورتر}
\begin{itemize}
\item شدیدترین موارد \lr{VUR}.
\item مثانه \lr{overdistended} می‌شود.
\item اکثراً در پسرها؛ می‌تواند یک‌طرفه یا دوطرفه باشد.
\end{itemize}

\VURpage{۲}{20261005\_024357191.jpg}
\section*{\lr{Primary VUR}}
\begin{itemize}
\item یک بیماری \lr{AD} با وریبل پنترنس.
\item ۳۵\% خواهر–برادرهای بیمار \lr{VUR}، \lr{VUR} دارند.
\item احتمال ابتلای خواهر–برادر ارتباطی به جنسیت کودک مبتلا یا گرید \VURunclear{ادامهٔ عبارت در حاشیهٔ چپ: ندارد}.
\end{itemize}
۱۲\% کودکان با \lr{VUR} بدون علامت، درجاتی از اسکار کلیه را دارند.

۵۰\% فرزندان مادر با \lr{VUR}، مبتلا به \lr{VUR} می‌شوند.

\section*{اندیکاسیون انجام \lr{VCUG} در خواهر–برادر فرد مبتلا}
یا سیستوگرام رادیونوکلئوتید.
\begin{itemize}
\item \lr{Febrile UTI} در فرد مبتلا به \lr{VUR}.
\item در سونوگرافی کلیهٔ خواهر–برادر:
\begin{itemize}
\item عدم قرینگی اندازهٔ دو کلیه.
\item وجود شواهد ابنورمالیتی کورتکس.
\end{itemize}
\item در غیر این صورت، غربالگری اختیاری است.
\end{itemize}
\subsection*{شواهد به نفع \lr{VUR} در سونوگرافی پیش از تولد}
\begin{itemize}
\item هیدرونفروز.
\item هیدرویورترونفروز.
\end{itemize}
\section*{گریدینگ \lr{VUR}}
\begin{description}
\item[\lr{I}:] میزنای دیلاته نیست. ادرار به سیستم جمع‌کنندهٔ فوقانی نمی‌رسد.
\item[\lr{II}:] میزنای دیلاته نیست. ادرار به سیستم جمع‌کنندهٔ فوقانی می‌رسد.
\item[\lr{III}:] میزنای دیلاته است و/یا فورنیکس‌های کالیس \VURunclear{واژهٔ پیش از «بلانت»: ریز} بلانت شده‌اند.
\item[\lr{IV}:] میزنای به‌طور \lr{gross} دیلاته شده است.
\item[\lr{V}:] تورتوزیته و از دست رفتن \VURunclear{واژهٔ زیر چاپ YASHA: فورنیکس؟}؛ ظاهر پاپیلاری نیز از بین رفته است.
\end{description}

\VURpage{۳}{20261005\_024401427.jpg}
\section*{\lr{VUR}}
\begin{itemize}
\item پرایمری: به‌طور مادرزادی اختلال عملکرد دریچه وجود داشته است.
\item پرایمری همراه با سایر مالفورماسیون‌های یورترووزیکال جانکشن.
\begin{itemize}
\item علل: دابلکیشن میزنای، یورتروسل + دابلکیشن، یورترال اکتوپی، دایورتیکول پارایورترال.
\end{itemize}
\item ثانویه به فشار اینتراوزیکال.
\begin{itemize}
\item علل: مثانهٔ نوروپاتیک، انسداد خروجی مثانه.
\end{itemize}
\item ثانویه به پروسه‌های التهابی.
\begin{itemize}
\item علل: سیستیت باکتریال شدید، جسم خارجی، سنگ مثانه.
\end{itemize}
\item ثانویه به جراحی.
\end{itemize}
۲۵\% کودکان با مثانهٔ نوروپاتیک (مثلاً در مننگومیلوسل)، \lr{VUR} پیدا می‌کنند.

\section*{تظاهرات بالینی}
\begin{itemize}
\item معمولاً طی بررسی‌های لازم برای \lr{UTI} به \lr{VUR} پی می‌بریم.
% In the image, UTI was written above a crossed-out VUR.
\item معمولاً ۸۰\% دختر و میانگین سنی ۲–۳ سال دارند.
\item مواردی که پیش از تولد هیدرونفروز دارند، ۸۰\% پسر هستند.
\item بدون علامت: صرفاً تب.
\item علامت‌دار: تب + درد شکم یا دیزوری.
\end{itemize}
\section*{تشخیص \lr{VUR}}
\begin{itemize}
\item نیازمند تعبیهٔ سوند ادراری.
\item عکس‌برداری حین پر شدن و خالی شدن مثانه.
\end{itemize}

\VURpage{۴}{20261005\_024408351.jpg}
\section*{تشخیص (ادامه)}
\begin{itemize}
\item \lr{VUR} حین پر شدن مثانه: \lr{low pressure VUR}.
\item \lr{VUR} حین تخلیهٔ مثانه: \lr{high pressure VUR}.
\end{itemize}
میزان اشعه‌ای که کودک از رادیونوکلئوتید سیستوگرام دریافت می‌کند خیلی کمتر از \lr{VCUG} است.

می‌توان قبل از پروسیجر به کودک میدازولام (\lr{nasal/oral}) یا پروپوفول داد.

\lr{Indirect cystography}: تنها ۷۵\% موارد \lr{VUR} را تشخیص می‌دهد.

پس از تأیید تشخیص \lr{VUR} باید کلیه‌ها را از نظر اسکار بررسی کنیم:
\begin{itemize}
\item سونوگرافی یا سینتی‌گرافی کلیه.
\end{itemize}
کودک \lr{VUR} باید از نظر \lr{BBD} نیز بررسی شود.

کودک \lr{VUR} باید قد، وزن و فشار خونش مانیتور شود.

اگر در بررسی کلیه شواهد اسکار دیده شد: \lr{Cr}، \lr{U/A}، پروتئین ادرار.

\section*{سیر بالینی}
\begin{itemize}
\item با افزایش سن و بلوغ، \lr{VUR} اغلب رفع می‌شود.
\item هرچه گرید پایین‌تر باشد، احتمال رفع خودبه‌خودی آن بیشتر است.
\item گرید \lr{I} و \lr{II} چه یک‌طرفه چه دوطرفه، در هر سنی معمولاً رفع می‌شوند.
\end{itemize}

\VURpage{۵}{20261005\_024415892.jpg}
\section*{سیر بالینی (ادامه)}
\begin{itemize}
\item گرید \lr{III} هرچه سن کمتر باشد و اگر یک‌طرفه باشد، با احتمال بالاتر رفع می‌شود.
\item گرید \lr{IV} دوطرفه، بعید است که رفع شود.
\item گرید \lr{V} معمولاً رفع نمی‌شود.
\end{itemize}
میانگین سن رفع \lr{VUR}، ۶ یا ۷ سالگی است.

مهم‌ترین \lr{RF} برای \lr{UTI} تب‌دار راجعه و اسکار کلیه:
\begin{itemize}
\item گرید \lr{III--V} باشد.
\item وجود \lr{BBD}.
\end{itemize}
\lr{VUR} استریل معمولاً آسیب کلیه ایجاد نمی‌کند مگر در موارد \lr{high-pressure}:
\begin{itemize}
\item \lr{posterior urethral valve}.
\item مثانهٔ نوروپاتیک نوروژنیک.
\item مثانهٔ \lr{non neurogenic}: سندرم هینمن (\lr{Hinman Syn})؛ \lr{non-psychiatric}.
\end{itemize}
\section*{درمان}
\begin{itemize}
\item هدف درمان: جلوگیری از پیلونفریت، آسیب کلیه و سایر عوارض.
\item درمان باید برای هر کسی به‌طور جداگانه و اختصاصی انتخاب شود.
\item دفع ادرار زمان‌بندی‌شده، اطمینان از عدم یبوست، تشخیص و درمان زودهنگام \lr{UTI}.
\item کدام بیماران با \lr{observation}: گرید \lr{I--II}، کلیهٔ \lr{NL}، بدون سابقهٔ پیلونفریت.
\end{itemize}

\VURpage{۶}{20261005\_024420188.jpg}
\section*{درمان آنتی‌میکروبیال}
\begin{itemize}
\item قبلاً پروفیلاکسی روزانه توصیه می‌شده.
\item در حال حاضر توصیه نمی‌شود.
\item ریسک \lr{UTI} راجعه در کدام بیماران بالاتر است؟
\begin{itemize}
\item گرید \lr{III} یا \lr{IV}.
\item \lr{BBD}.
\item بیماری که اولین \lr{UTI} او تب‌دار بوده است.
\end{itemize}
\item \lr{ABT} پروفیلاکسی پس از \lr{UTI}، احتمال \lr{UTI} بعدی را کم می‌کند ولی احتمال مقاومت باکتری بالا می‌رود.
\item به نظر نمی‌رسد پروفیلاکسی احتمال اسکار کلیه را خیلی کاهش دهد.
\end{itemize}
\section*{اندیکاسیون‌های پروفیلاکسی آنتی‌بیوتیکی}
\begin{itemize}
\item سن زیر ۱ سال.
\item وجود \lr{BBD} همراه با \lr{VUR}.
\item اگر سابقهٔ یک نوبت \lr{UTI} تب‌دار داشته باشد.
\end{itemize}
اندیکاسیون قطعی برای جراحی وجود ندارد.
